% !TEX program = xelatex
\documentclass[UTF8,12pt,a4paper]{ctexart}

\usepackage[a4paper,margin=2.5cm]{geometry}
\usepackage{amsmath,amssymb,bm}
\usepackage{booktabs}
\usepackage{graphicx}
\usepackage{float}
\usepackage{subcaption}
\usepackage{xcolor}
\usepackage{hyperref}
\usepackage{enumitem}

\hypersetup{
    colorlinks=true,
    linkcolor=blue!55!black,
    citecolor=blue!55!black,
    urlcolor=blue!55!black
}

\newcommand{\ReportTitle}{一维 Riemann 问题的 MacCormack 数值结果与讨论}
\newcommand{\ReportAuthor}{填写姓名}
\newcommand{\ReportCourse}{填写课程名称}
\newcommand{\ReportDate}{\today}

\newcommand{\figureplaceholder}[2]{
\begin{figure}[H]
    \centering
    \fbox{
        \parbox[c][5cm][c]{0.82\textwidth}{
            \centering
            将图片放入 \texttt{figures/} 后，\\
            用 \texttt{\textbackslash includegraphics} 替换此占位框。
        }
    }
    \caption{#1}
    \label{#2}
\end{figure}
}

\title{\ReportTitle}
\author{\ReportAuthor\\\ReportCourse}
\date{\ReportDate}

\begin{document}

\maketitle
\begin{abstract}
本文展示 MacCormack 格式求解一维 Riemann 问题的数值结果，并与
Project 0 的半解析结果进行比较。此处填写主要算例、参数研究和最终结论。

\textbf{关键词：} 数值结果；Sod 激波管；人工粘性；CFL；网格收敛性
\end{abstract}

\tableofcontents
\newpage

\section{算例与计算参数}

\subsection{测试算例}

列出本报告实际采用的 Riemann 初始条件。

\begin{table}[H]
    \centering
    \caption{Riemann 问题初始条件}
    \label{tab:cases}
    \begin{tabular}{lcccccc}
        \toprule
        算例 & $\rho_L$ & $u_L$ & $p_L$ & $\rho_R$ & $u_R$ & $p_R$\\
        \midrule
        Sod & 1.000 & 0.000 & 1.000 & 0.125 & 0.000 & 0.100\\
        Lax & 0.445 & 0.698 & 3.528 & 0.500 & 0.000 & 0.571\\
        自定义 & -- & -- & -- & -- & -- & --\\
        \bottomrule
    \end{tabular}
\end{table}

\subsection{基准计算参数}

\begin{table}[H]
    \centering
    \caption{基准数值参数}
    \label{tab:parameters}
    \begin{tabular}{lc}
        \toprule
        参数 & 数值\\
        \midrule
        计算域 $[x_{\min},x_{\max}]$ & $[0,1]$\\
        初始间断位置 $x_0$ & $0.5$\\
        网格数 $N_x$ & 501\\
        CFL & 0.5\\
        比热比 $\gamma$ & 1.4\\
        终止时间 $t_{\max}$ & 0.2\\
        人工粘性系数 & 填写最终数值\\
        \bottomrule
    \end{tabular}
\end{table}

\subsection{评价指标}

说明将比较哪些量：

\begin{itemize}[leftmargin=2em]
    \item 密度 $\rho$；
    \item 速度 $u$；
    \item 压力 $p$；
    \item 激波、接触间断和稀疏波的位置；
    \item 与半解析解的误差；
    \item 振荡幅度和间断厚度。
\end{itemize}

\section{基准算例结果}

\subsection{Sod 激波管}

描述稀疏波、接触间断和激波的传播方向与数值结果。

\figureplaceholder{Sod 算例密度分布}{fig:sod-density}
\figureplaceholder{Sod 算例速度分布}{fig:sod-velocity}
\figureplaceholder{Sod 算例压力分布}{fig:sod-pressure}

\subsection{Lax 问题}

填写 Lax 算例的结果和主要波结构。

\figureplaceholder{Lax 算例数值结果}{fig:lax-result}

\subsection{其他 Riemann 算例}

根据最终完成情况增加强膨胀、双激波或接触间断算例。每个算例至少记录：

\begin{itemize}[leftmargin=2em]
    \item 初始条件；
    \item 网格数、CFL、终止时间和人工粘性系数；
    \item 是否出现负压、负密度或非物理振荡；
    \item 数值格式能否稳定完成计算。
\end{itemize}

\section{与半解析解的比较}

\subsection{对比方法}

说明 Project 0 输出与 MacCormack 输出如何对齐，例如统一采样坐标或插值。

\subsection{曲线对比}

\figureplaceholder{数值解与半解析解的密度对比}{fig:compare-density}
\figureplaceholder{数值解与半解析解的速度对比}{fig:compare-velocity}
\figureplaceholder{数值解与半解析解的压力对比}{fig:compare-pressure}

\subsection{误差计算}

可以使用离散 $L_1$、$L_2$ 和 $L_\infty$ 误差：

\begin{align}
L_1 &= \frac{1}{N_x}\sum_{i=1}^{N_x}
\left|q_i^{\mathrm{num}}-q_i^{\mathrm{ref}}\right|,\\
L_2 &= \sqrt{\frac{1}{N_x}\sum_{i=1}^{N_x}
\left(q_i^{\mathrm{num}}-q_i^{\mathrm{ref}}\right)^2},\\
L_\infty &= \max_i
\left|q_i^{\mathrm{num}}-q_i^{\mathrm{ref}}\right|.
\end{align}

\begin{table}[H]
    \centering
    \caption{数值解误差}
    \label{tab:error}
    \begin{tabular}{lccc}
        \toprule
        物理量 & $L_1$ & $L_2$ & $L_\infty$\\
        \midrule
        密度 $\rho$ & -- & -- & --\\
        速度 $u$ & -- & -- & --\\
        压力 $p$ & -- & -- & --\\
        \bottomrule
    \end{tabular}
\end{table}

\section{参数影响与结果讨论}

\subsection{人工粘性的影响}

对比关闭和开启人工粘性的结果，讨论：

\begin{itemize}[leftmargin=2em]
    \item 间断附近振荡是否减弱；
    \item 激波是否被过度抹宽；
    \item 接触间断的分辨率变化；
    \item 不同人工粘性系数的影响。
\end{itemize}

\figureplaceholder{不同人工粘性系数的结果比较}{fig:viscosity}

\subsection{CFL 的影响}

建议至少比较三组满足稳定条件的 CFL，例如 0.2、0.5 和 0.8。

\begin{table}[H]
    \centering
    \caption{不同 CFL 的计算表现}
    \label{tab:cfl}
    \begin{tabular}{cccc}
        \toprule
        CFL & 时间步数 & 是否稳定 & 主要现象\\
        \midrule
        0.2 & -- & -- & --\\
        0.5 & -- & -- & --\\
        0.8 & -- & -- & --\\
        \bottomrule
    \end{tabular}
\end{table}

\subsection{网格分辨率的影响}

比较不同 $N_x$，讨论激波厚度、误差和计算成本。

\figureplaceholder{不同网格分辨率的结果比较}{fig:grid}

\subsection{数值异常与适用范围}

记录程序运行中观察到的负压、负密度、NaN、步数上限或边界影响。
区分以下两类结论：

\begin{itemize}[leftmargin=2em]
    \item 格式本身的限制；
    \item 当前程序实现仍需改进的部分。
\end{itemize}

\section{结论}

建议按以下结构填写：

\begin{enumerate}[leftmargin=2em]
    \item MacCormack 格式是否能够捕捉一维 Riemann 问题的基本波结构；
    \item 与半解析解相比，主要误差出现在哪里；
    \item 人工粘性对振荡和间断分辨率的影响；
    \item CFL 与网格分辨率对结果的影响；
    \item 当前方法的适用范围和下一步改进方向。
\end{enumerate}

\section*{参考资料}

\begin{enumerate}[leftmargin=2em]
    \item 填写课程教材或讲义。
    \item 填写 Project 0 半解析解使用的理论资料。
    \item 填写 MacCormack 格式及人工粘性相关资料。
\end{enumerate}


\end{document}
